\chapter{Lanzador de dados}
\label{chap:dados}
Esta es la función de dados de la aplicación.
Permite elegir el tipo de dado y cuántos se lanzan, entre 1 y 20, y muestra el resultado de cada uno y la suma.

Se sigue la misma paleta de colores que en el resto de la aplicación.

\begin{figure}[H]
  \centering
  \includegraphics[width=0.4\textwidth]{dados}
\end{figure}

\section{Enumeración TipoDado}
Cada tipo de dado se define una sola vez con su número de caras.

\begin{lstlisting}
enum class TipoDado(val caras: Int) {
    D4(4),
    D6(6),
    D8(8),
    D10(10),
    D12(12),
    D20(20),
    D100(100);

    override fun toString() = "d$caras"
}
\end{lstlisting}

La función \texttt{toString()} está sobrescrita para que cada dado se muestre como d4, d6, etc.
El desplegable usa ese texto directamente, así que se le puede pasar la lista de valores del enum sin convertirla antes.
Empieza seleccionado el dado de 6 caras.

Añadir un dado nuevo es añadir una línea.

\section{Cambiar la cantidad de dados}
Los botones de más y menos llaman a la función \texttt{cambiarCantidad}, que recibe un $+1$ o un $-1$.

\begin{lstlisting}
private fun cambiarCantidad(cambio: Int) {
    cantidad = (cantidad + cambio).coerceIn(minDados, maxDados)
    textoCantidad.text = cantidad.toString()
}
\end{lstlisting}

Con \texttt{coerceIn} la cantidad nunca baja del mínimo ni pasa del máximo, así que no hacen falta más comprobaciones.

\section{Lanzar los dados}
La función \texttt{lanzar} mira qué dado está elegido y genera un número aleatorio por cada dado.

\begin{lstlisting}
val tipo = TipoDado.entries[spinnerTipo.selectedItemPosition]
val resultados = List(cantidad) { Random.nextInt(1, tipo.caras + 1) }
\end{lstlisting}

Se suma 1 a las caras porque \texttt{Random.nextInt} no incluye el último número.

\subsection{Dibujar los resultados}
Primero se borran los dados de la tirada anterior con \texttt{removeAllViews()}.
Después se crea un cuadrado por cada resultado y cada 5 dados se crea una nueva fila.

\begin{lstlisting}
if (i % 5 == 0) {
    fila = LinearLayout(this)
    fila.orientation = LinearLayout.HORIZONTAL
    layoutDados.addView(fila)
}
\end{lstlisting}

Los cuadrados son \texttt{TextView} creados desde el código, porque no se sabe cuántos hacen falta hasta que se lanza.
Por eso el color, la letra y el tamaño también se ponen ahí y no en el XML.
Las medidas se multiplican por la densidad de la pantalla para pasarlas de dp a píxeles.

Al final se muestra la suma de todos los dados.